\chapter{Literature Review}

\section{Introduction}

The increasing adoption of machine learning models in sensitive domains such as finance, healthcare, and authentication has raised concerns regarding their verifiability, privacy, and correctness. Zero-Knowledge Proofs (ZKPs) have emerged as a powerful cryptographic tool to validate computations without revealing sensitive model parameters or data. This chapter presents a structured review of recent literature on verifiable machine learning, efficient ZKP techniques, and ZKP-based optimization strategies, focusing on contributions relevant to building verifiable logistic regression models using ZKPs.

\section{Scalable Zero-Knowledge Proofs for Non-Linear Functions in Machine Learning}

Hao et al. (2025) introduced a scalable ZKP framework for handling non-linear functions in machine learning, including sigmoid, ReLU, normalization, and other activation functions~\cite{cryptoeprint:2025/507}.  
Their framework employs digit decomposition and lookup-table-based computation to achieve 50 to 179 $\times$ runtime improvement compared to previous ZKP systems.

\textbf{Insights:}  
This work underscores the effectiveness of lookup-table-based techniques for efficient non-linear function evaluation, directly influencing our approach to implementing sigmoid computation in logistic regression circuits.

\textbf{Limitations:}  
The system primarily targets deep learning models, with limited optimizations specific to logistic regression.

\section{Mystique: Efficient Conversions for Zero-Knowledge Proofs}

Weng et al. (2021) developed Mystique, a ZKP framework enabling efficient conversions between arithmetic and Boolean circuits to support large-scale machine learning inference~\cite{cryptoeprint:2021/730}.  
The system enhances interoperability and performance for matrix-based operations in ZKP environments.

\textbf{Insights:}  
Mystique demonstrates the benefits of modular circuit construction and efficient arithmetic–Boolean translation, which inspire the layered design of our ZKLR circuit.

\textbf{Limitations:}  
The system is optimized for complex neural networks, making it heavy for lightweight models such as logistic regression.

\section{Zero-Knowledge Proofs for Decision Tree Predictions and Accuracy}

Zhang et al. (2020) proposed a ZKP protocol for verifying decision-tree predictions and accuracy claims without revealing the internal structure of the model~\cite{10.1145/3372297.3417278}.  
The system encodes branching logic into arithmetic constraints, enabling verifiable ML predictions.

\textbf{Insights:}  
This approach motivates our system’s accuracy-verification mechanism for logistic regression predictions.

\textbf{Limitations:}  
Decision-tree structures are discrete and simpler to encode, while continuous sigmoid-based functions require more complex circuit construction.

\section{Proof of Training for Logistic Regression}

Garg et al. (2023) introduced a proof-of-training (zkPoT) system for logistic regression that uses a hybrid SNARK and MPC-in-the-head design to ensure training correctness over large datasets~\cite{cryptoeprint:2023/1345}.

\textbf{Insights:}  
Their work demonstrates how logistic regression training can be verifiably proven in zero-knowledge, laying conceptual groundwork for inference verification.

\textbf{Limitations:}  
The system focuses primarily on the training phase rather than inference or accuracy validation.

\section{Efficient Zero-Knowledge Proofs for Deep Learning Training}

Sun et al. (2023) presented zkDL, an efficient ZKP framework for verifying deep learning training processes~\cite{cryptoeprint:2023/1174}.  
They introduced zkReLU and structured circuit optimizations to improve proving performance.

\textbf{Insights:}  
zkDL’s lookup-based and batch-proof aggregation strategies inspired similar design patterns for efficient circuit implementation in our system.

\textbf{Limitations:}  
The framework is specialized for deep learning, not directly optimized for lightweight logistic regression models.

\section{Benchmarking Non-Interactive Zero-Knowledge Proof Protocols}

El-Hajj and Roelink (2024) conducted a benchmark comparison of zk-SNARK, zk-STARK, and Bulletproof systems, analyzing performance metrics such as proof size, verification time, and prover cost~\cite{info15080463}.

\textbf{Insights:}  
Their findings justify our adoption of PLONK for its balanced trade-off between succinct proof size and computational efficiency.

\textbf{Limitations:}  
Their analysis is based on hash-based circuits and does not address ML-specific computational constraints.

\section{ZKP-Based Verifiable Machine Learning Survey}

Xing et al. (2025) provided a comprehensive survey of ZKP-based verifiable decentralized machine learning frameworks~\cite{article}.  
The survey categorizes ZKP techniques and identifies challenges in building trustworthy privacy-preserving ML systems.

\textbf{Insights:}  
This survey offers a broad contextual understanding, helping align our architectural choices with emerging ZKML design trends.

\textbf{Limitations:}  
It remains theoretical and does not detail implementation strategies for logistic regression under ZKP constraints.

\section{Research Gaps Identified}

Despite significant progress in zero-knowledge machine learning, several research gaps persist that motivate the development of our Zero-Knowledge Logistic Regression (ZKLR) framework.  
First, there is a limited focus on lightweight machine learning models such as logistic regression. Most existing ZKP systems are designed for complex deep learning architectures, overlooking simpler yet practically relevant models~\cite{cryptoeprint:2023/1174}.  
Second, scalability challenges remain unresolved, as many frameworks have not been evaluated across datasets of varying scales, restricting their applicability to real-world environments where dataset sizes fluctuate~\cite{cryptoeprint:2021/730}.  
Third, the computation of non-linear functions—particularly the sigmoid function—continues to pose a bottleneck. Efficient and ZKP-friendly sigmoid approximations remain an open problem in circuit design~\cite{cryptoeprint:2025/507}.  
Fourth, existing works generally adopt fragmented verification strategies, validating either individual prediction correctness or aggregate model accuracy, but rarely providing an integrated framework that addresses both concerns~\cite{10.1145/3372297.3417278}. 
Finally, dataset adaptivity remains underexplored. Current systems lack flexibility to handle datasets with varying sizes and feature dimensions, limiting their deployment across diverse machine learning applications~\cite{article}.

These gaps collectively underscore the need for an efficient, scalable, and adaptable ZKLR system capable of performing verifiable logistic regression inference while maintaining strong privacy and performance guarantees~\cite{kates2020plonk}.

\section{Conclusion}

The reviewed literature establishes foundational approaches for applying zero-knowledge proofs to machine learning verification. Our work addresses the identified limitations through an integrated system combining modular circuit architecture, lookup-table-based sigmoid approximation, and recursive proof composition for scalable accuracy verification. These contributions extend the state of the art in verifiable and privacy-preserving logistic regression inference.

\clearpage
